\section*{ID9151: Relation: < 0). the Radial Acceleration Is a\_r =}
\paragraph{Metadata.}
PiID: 6959469955657 | RTEID: RTE:P35535.3652:T2.0027 | UUID: b4b636b6-6a30-5033-b0cd-5661969e3d51

\paragraph{Canonical Statement.}
Derivable Consequences: Amplified radial pressure gradients; crosssymmetry scaling laws; design of vortex generators; improved vortex strength. Falsifiability / Test Handle: Compare flow fields with and without the cross; measure symmetry and inflow rates. 1.0.24 ID 123: Radial Pressure Gradient and Acceleration Field Type: Definition Formal Statement: From Bernoulli’s equation P + ¡ \$ρ\$ vš = constant, increased velocity at the centre hole implies a pressure decrease toward the centre \$(∂P/∂r\$ < 0). The radial acceleration is a\_r = \$–(1/ρ) ∂P/∂r.\$ In a cross geometry, forces from four directions sum, amplifying the pull by a factor roughly proportional to 2 compared with a simple circular opening. Interpretive Meaning: This describes how velocity gradients translate into pressure gradients and acceleration toward the centre; the cross symmetry magnifies the effect. Dependencies: 122 Derivable Consequences: Scaling of pull str

\paragraph{Canonical Equation.}
\[
< 0). The radial acceleration is a_r =
\]

\paragraph{Dictionary.}
Derivable Consequences: Amplified radial pressure gradients; crosssymmetry scaling laws; design of vortex generators; improved vortex strength.

\paragraph{Encyclopedia.}
Relation: < 0). the Radial Acceleration Is a\_r = is an inequality / bound, written < 0). The radial acceleration is a\_r =. It states an inequality or bound that constrains the admissible range of the quantity. Within the Trawin Zero Point registry it serves as one element of the framework's mathematical narrative.
